\documentclass[a4paper,11pt,openany]{report}
\usepackage{../preamble/agenticboost}

\renewcommand{\sprintId}{S-001.5R}

\title{%
  \textsc{AgenticCareerBoost}\\[0.4em]
  \Large Sprint S-001.5R: System Review and Publication Governance\\[0.3em]
  \large Missing S-001.5 Documentation $\cdot$ Failure Analysis $\cdot$ Next Controls%
}
\author{Didac Llorens \\ \small \href{https://github.com/DidacLL}{github.com/DidacLL}}
\date{April 2026 --- v0.1}

\begin{document}
\maketitle
\tableofcontents
\newpage

% ============================================================================
\chapter{Executive Summary}
\label{ch:executive}
% ============================================================================

\takeaway{S-001.5 achieved the urgent public-profile objective: a recruiter
landing surface, CV, GitHub cleanup package, LinkedIn/JobTeaser copy, outreach
drafts, and event target shortlist were produced. The review also exposed a
governance problem: publication, site architecture, generated artifacts, and
social planning were coupled too tightly for the system to remain clean under
deadline pressure.}

S-001.5 was the UOC fair readiness sprint. Its purpose was practical rather
than architectural: make the public profile coherent enough for recruiter
inspection before the VI Fira Virtual d'Ocupacio de la UOC. The sprint closed
with local artifacts and green build checks, but GitHub Pages still required a
repository setting change outside the workspace.

S-001.5R is the system review that documents what happened after that closure.
It fills the missing formal documentation gap for S-001.5 and records the
flaws found during the human review:

\begin{itemize}
  \item direct commits reached \fname{main} without an enforced review gate;
  \item the site existed in more than one implementation shape;
  \item generated files appeared in source-oriented folders;
  \item the Pages publishing model and repository settings did not match;
  \item sprint execution and social publishing plans became coupled;
  \item public-tone safeguards had to be repeatedly rechecked.
\end{itemize}

The corrective direction is not to erase S-001.5. It delivered useful public
material. The correction is to separate source, build output, publication, and
social distribution so future sprints can move quickly without blurring the
truth hierarchy.

\pitfall{The mistake would be treating S-001.5 as a failed sprint. It was an
effective deadline sprint with insufficient publication governance. The
artifact value should be kept; the workflow leak should be fixed.}


% ============================================================================
\chapter{What S-001.5 Produced}
\label{ch:s0015}
% ============================================================================

\takeaway{S-001.5 converted the S-001 positioning work into concrete public
assets: profile copy, event outreach, a fair-ready site, a CV, and a checklist
for user-controlled account actions.}

\section{Sprint Goal}

The sprint goal was to make the public profile event-ready by applying the
S-001 positioning across GitHub, the public site, LinkedIn/JobTeaser copy,
outreach messaging, and target-company triage. The core positioning remained:
ML/Data Engineering as the primary lane, Agentic Systems as the differentiator,
and documented systems work as the proof surface.

\section{Completed Artifact Set}

\begin{table}[htbp]
\centering
\caption{S-001.5 artifact map.}
\label{tab:s0015-artifacts}
\small
\begin{tabularx}{\textwidth}{l X X}
\toprule
\textbf{Area} & \textbf{Artifact} & \textbf{Purpose} \\
\midrule
Research & \pathref{state/research/s0015-uoc-fair-targets.md} & Rank UOC fair targets by fit, message angle, and risk. \\
GitHub profile & \pathref{content/social/drafts/2026-04-s0015-github-cleanup.md} & Prepare README, bio, repo descriptions, topics, pin order, and archive guidance. \\
Site & \pathref{site/} & Provide a 90-second recruiter scan with JSON-LD, project evidence, CV link, and contact route. \\
CV & \pathref{content/reports/tex/guides/didac-llorens-cv.tex} & Produce a formal CV aligned with the ML/Data plus Agentic Systems posture. \\
LinkedIn/JobTeaser & \pathref{content/social/drafts/2026-04-s0015-linkedin-jobteaser-kit.md} & Supply paste-ready profile copy without unsupported credentials. \\
Outreach & \pathref{content/social/drafts/2026-04-s0015-uoc-fair-outreach.md} & Draft event intros, interview requests, follow-ups, and company-specific angles. \\
Execution & \pathref{state/summaries/s0015-execution-checklist.md} & Separate local repository work from user-controlled external account actions. \\
\bottomrule
\end{tabularx}
\end{table}

\section{Review Status}

PairCheck passed the public copy and recruiter site review. The checks confirmed
that the copy avoided the legacy site, unsupported credentials, fabricated
metrics, and apology framing. The site review confirmed the recruiter scan,
project evidence, curriculum page, and CV link. The remaining blocker was not
content quality. It was publication state: the build path had to be made
consistent with GitHub Pages settings.


% ============================================================================
\chapter{Detected Flaws}
\label{ch:flaws}
% ============================================================================

\takeaway{The failures were system-boundary failures. Different agents solved
urgent local problems, but the repository did not yet enforce the boundaries
between reviewed source, generated output, deployable site artifact, and public
campaign planning.}

\section{Direct-Main Commit}

The git history shows S-001.5 closure and publication commits present directly
on \fname{main}, including \fname{a1db91d}, \fname{8642f6d}, \fname{10bb463},
\fname{e812508}, \fname{f601700}, and \fname{4abef07}. That is evidence of a
missing governance gate, not evidence that those commits were useless.

The practical risk is simple: a sprint can be declared reviewed after the fact
even when the repository never forced a review-before-merge path. For a project
whose public claim is disciplined agentic workflow design, that weakens the
claim. S-001.5R therefore treats branch protection and one aggregate required
CI check as publication infrastructure, not administrative polish.

\section{Duplicate Site Implementations}

The repository accumulated multiple site shapes:
\pathref{site/starter/} as the original Jekyll scaffold,
\pathref{content/site/} as Markdown content intended for the site, and
\pathref{site/public/} as a static deadline artifact. The static path was a
valid emergency adaptation because the fair deadline was immediate and the
Pages Actions setup failed. It still left the next sprint with a design
question: which tree is canonical?

The answer should be explicit. S-002 should keep one site source model and one
deployable output model. Until then, documentation must say when a file is
source, generated mirror, or emergency static artifact.

\section{Generated Files in Source Folders}

Generated PDFs appeared under source-oriented TeX directories, while the
intended publication folder is \pathref{content/reports/build/} and local TeX
output belongs under \pathref{content/reports/tex/build/}. This matters because
source review becomes noisy when build artifacts sit beside editable report
sources.

The right invariant is already present in the LaTeX infrastructure: source
files live under \pathref{content/reports/tex/}; local build output goes to
\fname{content/reports/tex/build/}; promoted public PDFs go to
\pathref{content/reports/build/}. S-001.5R should preserve that separation.

\section{Pages Publishing Mismatch}

The review logs record two different Pages models. The original workflow used
GitHub Pages Actions, but \fname{actions/configure-pages} could not create or
configure the Pages site with the default token. The deadline workaround
published static HTML to \fname{gh-pages}. After that, the repository still
returned 404 until the human enabled Pages with the matching source setting.

The flaw was not GitHub Pages itself. The flaw was a mismatch between local CI
design, remote repository settings, and sprint closure language. A green build
does not mean a live site unless the configured Pages source matches the
workflow's publishing model.

\section{Sprint/Social Coupling}

S-001.5 mixed sprint execution with public-distribution readiness: GitHub
cleanup, LinkedIn/JobTeaser copy, outreach drafts, social plan updates, and
site deployment all moved in the same deadline window. That was useful for the
fair, but it made state synchronization fragile. The ContentSync review found
that the social plan still sequenced only S-000 posts until it was corrected.

The rule for future sprints should be stricter: sprint artifacts can be created
inside the sprint, but publication scheduling should remain gated by verified
artifact availability, live links, and CommunityManager review.

\section{Public-Tone Drift}

The brand allows a controlled edge in public-facing artifacts, but the review
history shows recurring tone risks: hype-adjacent language, excess sarcasm,
unsupported statistics, and junior-coded legacy copy. S-001.5 public copy
passed the final check, but only because PairCheck explicitly revalidated the
tone and evidence constraints.

The system should treat tone as a testable public-safety condition. Public copy
must remain technical, direct, evidence-backed, and free of generic AI hype,
founder-LARP framing, apology framing, and unsupported metrics.


% ============================================================================
\chapter{What S-001.5R Changed}
\label{ch:changes}
% ============================================================================

\takeaway{S-001.5R does not rewrite the sprint artifacts. It records the review
findings and defines the controls required before the next public push is
treated as complete.}

The companion governance log, \pathref{state/logs/S-001.5R/ci-pages-governance.md},
records the local CI/Pages changes under review:

\begin{itemize}
  \item a stable aggregate workflow named \fname{required-ci};
  \item validation for pytest, markdown lint, internal links, static site root,
    and LaTeX compilation;
  \item a shift away from direct \fname{gh-pages} branch commits and toward
    GitHub Pages Actions that build Jekyll from \pathref{site/} and deploy
    \fname{site/\_site/}.
\end{itemize}

Those changes still require remote repository settings that cannot be applied
from the local workspace:

\begin{enumerate}
  \item enable Pages with source \textbf{GitHub Actions};
  \item protect \fname{main} or add an equivalent ruleset;
  \item require the aggregate \fname{required-ci} status before merge;
  \item remove fragmented required checks after the aggregate check has run.
\end{enumerate}

\pitfall{A required check is useful only if it is required by repository
settings. A workflow file in the repo is evidence of intent; branch protection
is the enforcement mechanism.}


% ============================================================================
\chapter{Next Steps}
\label{ch:next}
% ============================================================================

\takeaway{The next sprint should not add more surfaces before stabilizing the
ones that already exist. S-002 should choose the canonical site architecture,
clean generated artifacts out of source paths, and keep publication gated by
live verification.}

\section{Immediate Human Actions}

\begin{enumerate}
  \item Set repository Pages to the model selected by S-001.5R. If the workflow
    uses GitHub Pages Actions, the remote setting must also use GitHub Actions.
  \item Add branch protection or a repository ruleset for \fname{main}.
  \item Require the aggregate \fname{required-ci} check once it has run on the
    default branch.
  \item Keep external profile changes user-controlled: LinkedIn, JobTeaser,
    GitHub profile settings, and recruiter outreach should not be automated.
\end{enumerate}

\section{S-002 Constraints}

\begin{itemize}
  \item Keep \pathref{site/} as the canonical site source tree.
  \item Keep generated PDFs out of editable TeX source folders.
  \item Keep \fname{site/\_site/} as generated build output, not checked-in
    source.
  \item Do not mark the public site complete until the URL has been verified
    live.
  \item Keep social publication blocked until the linked artifacts exist and
    tone checks pass.
\end{itemize}

\section{Closure Standard}

Future sprint closure should distinguish four states:

\begin{description}[style=nextline]
  \item[Source complete]
    The repo contains the intended source artifact.
  \item[Build green]
    CI has built and validated the artifact.
  \item[Published]
    The public surface is live at the intended URL.
  \item[Distributed]
    Social or external profile copy has been applied by the user.
\end{description}

S-001.5 reached source complete and build green. It did not fully reach
published until the Pages source setting matched the workflow. It did not reach
distributed automatically, by design, because external profile changes remain
human-controlled.

\end{document}
